\chapter{Data and Distributions}
\label{ch:data-distributions}

Here are twenty response times, in milliseconds, from one API endpoint over one minute:
\[
\begin{array}{cccccccccc}
12 & 14 & 15 & 15 & 17 & 18 & 18 & 19 & 21 & 22\\
23 & 23 & 24 & 26 & 28 & 31 & 35 & 42 & 58 & 140
\end{array}
\]
Is this endpoint fast? Is $140$ a problem or just bad luck? Would tomorrow look the same?
Staring at the list does not answer any of these. Twenty numbers are already too many to hold
in your head, and production systems log millions per hour. This chapter is about the first
move every statistician makes: turn a pile of numbers into a \emph{picture of how the values are
spread out}. That picture is called a \emph{distribution}, and nearly everything in the rest of
the book is a question about one.

\section{What statistics is for}
\label{sec:data-what}

Statistics does three jobs. Each one asks something different of the same data.

\begin{itemize}
  \item \textbf{Describe.} Summarize the data you have. ``Of these 20 requests, 15 finished in
  under 30 ms.'' No guessing is involved; it is bookkeeping done well.
  \item \textbf{Infer.} Use the data you have to say something about data you do not have.
  ``Across all of today's traffic, about 75\% of requests finish in under 30 ms, give or take
  a few points.'' Now there is uncertainty, and we must quantify it.
  \item \textbf{Predict.} Say what the next observation will look like. ``The next request will
  probably take between 12 and 60 ms.'' This is the job a machine learning model does.
\end{itemize}

\begin{figure}[H]
\centering
\begin{tikzpicture}[
  box/.style={draw, rounded corners, thick, align=center, minimum width=3.3cm, minimum height=1.5cm, font=\small},
  arr/.style={-Stealth, thick}
]
\node[box, fill=histfill!30] (d) at (0,0) {\textbf{Describe}\\what is in\\this sample};
\node[box, fill=targetgreen!20] (i) at (4.2,0) {\textbf{Infer}\\what is true of\\the population};
\node[box, fill=quotientorange!25] (p) at (8.4,0) {\textbf{Predict}\\what the next\\value will be};
\draw[arr] (d) -- (i);
\draw[arr] (i) -- (p);
\node[font=\footnotesize, align=center, below=0.25cm of d] {Chapters 1--4};
\node[font=\footnotesize, align=center, below=0.25cm of i] {Chapters 5--11};
\node[font=\footnotesize, align=center, below=0.25cm of p] {Chapters 12--15};
\end{tikzpicture}
\caption{The three jobs of statistics. Each builds on the one before: you cannot infer
anything about a population until you can describe a sample, and prediction is inference
about a single future value.}
\label{fig:data-three-jobs}
\end{figure}

This chapter and the next three are about describing. Description sounds like the easy part,
but many mistakes in practice happen here: the wrong chart, the wrong summary, a
misleading bin width, an outlier nobody looked at.

\section{Units, variables, and types}
\label{sec:data-variables}

\subsection{The shape of a data set}

A data set is a table. Each row is one \textbf{observational unit}: the thing you measured. Each
column is a \textbf{variable}: one property you recorded about every unit. Here is a slice of a
request log.

\begin{table}[H]
\centering
\small
\begin{tabular}{cllrcl}
\toprule
\texttt{request\_id} & \texttt{endpoint} & \texttt{status} & \texttt{latency\_ms} & \texttt{retries} & \texttt{priority}\\
\midrule
1 & \texttt{/search}  & 200 & 21.4  & 0 & high\\
2 & \texttt{/search}  & 200 & 18.9  & 0 & medium\\
3 & \texttt{/cart}    & 404 & 12.1  & 0 & low\\
4 & \texttt{/checkout}& 200 & 57.6  & 1 & high\\
5 & \texttt{/search}  & 503 & 140.2 & 3 & medium\\
\bottomrule
\end{tabular}
\caption{Five rows of a request log. Each row is one request (the unit). Each column is a variable.}
\label{tab:data-log}
\end{table}

Picking the unit is a real decision. If each row were a \emph{user} instead of a request, the
latency column would have to become something like ``average latency of that user's
requests,'' and the questions you could answer would change. Before any analysis, say out
loud: one row is one \underline{\hspace{1.5cm}}.

\subsection{Types of variables}

The type of a variable decides which summaries and charts make sense. There are two big
families.

\begin{definition}[Categorical and numerical variables]
A \textbf{categorical} variable places each unit into one of a set of groups (\texttt{endpoint},
\texttt{status}). A \textbf{numerical} variable records a quantity, a number for which
arithmetic means something (\texttt{latency\_ms}, \texttt{retries}).
\end{definition}

The test is simple: does averaging make sense? The average latency of requests 1 and 2 is
$20.15$ ms, a meaningful number. The ``average status code'' of a 200 and a 404 is 302, which is
a real HTTP code (a redirect) that neither request returned. Status codes are labels written
with digits. They are categorical.

Each family splits again.

\begin{itemize}
  \item A categorical variable is \textbf{nominal} if its groups have no natural order
  (\texttt{endpoint}: \texttt{/cart} is not ``more'' than \texttt{/search}), and
  \textbf{ordinal} if they do (\texttt{priority}: low $<$ medium $<$ high). Ordinal means
  you can rank the groups but the gaps between them have no fixed size: ``high minus medium''
  is not a number.
  \item A numerical variable is \textbf{discrete} if it takes separated values you can list,
  usually counts (\texttt{retries} is 0, 1, 2, \dots; never 1.5), and \textbf{continuous} if
  it can take any value in an interval (\texttt{latency\_ms} could be $21.4$ or $21.41$ or
  $21.4137$, limited only by the clock's precision).
\end{itemize}

\begin{figure}[H]
\centering
\begin{tikzpicture}[
  level distance=1.55cm,
  level 1/.style={sibling distance=6.2cm},
  level 2/.style={sibling distance=3.1cm},
  every node/.style={draw, rounded corners, align=center, font=\small, thick, minimum width=2.4cm, inner sep=4pt},
  edge from parent/.style={draw, -Stealth, thick}
]
\node[fill=gray!10] {Variable}
  child {node[fill=structurepurple!15] {Categorical\\(labels)}
    child {node[fill=structurepurple!8] {Nominal\\[-1pt]{\footnotesize endpoint, status}}}
    child {node[fill=structurepurple!8] {Ordinal\\[-1pt]{\footnotesize priority, rating}}}
  }
  child {node[fill=sourceblue!15] {Numerical\\(quantities)}
    child {node[fill=sourceblue!8] {Discrete\\[-1pt]{\footnotesize retries, heads}}}
    child {node[fill=sourceblue!8] {Continuous\\[-1pt]{\footnotesize latency, height}}}
  };
\end{tikzpicture}
\caption{The four variable types. The first split asks ``does arithmetic make sense?'' The
second asks ``is there an order?'' (categorical) or ``can I list the possible values?''
(numerical).}
\label{fig:data-types-tree}
\end{figure}

\begin{trapbox}
Numbers in a column do not make a variable numerical. Zip codes, user IDs, status codes, and
product SKUs are categorical. A plotting library will happily draw a histogram of zip codes;
the library did what it was told, and the mistake is yours. Going the other way, a 1--5 star rating is ordinal: treating
it as numerical (averaging stars) is common and often useful, but it assumes the step from 1 to
2 stars equals the step from 4 to 5, which nobody checked.
\end{trapbox}

\begin{recoverybox}
Classify each: (a) number of items in a shopping cart; (b) browser name; (c) time to first
byte; (d) T-shirt size S/M/L/XL. \emph{Answers:} discrete, nominal, continuous, ordinal.
\end{recoverybox}

\subsection{Populations and samples}

The twenty latencies above are not all the requests this endpoint will ever serve. They are
a \textbf{sample}: the units we actually observed. The \textbf{population} is the full set of
units we care about, such as every request this endpoint served today. A number that describes
the population (the true fraction of today's requests under 30 ms) is a \textbf{parameter}. The
same number computed from the sample is a \textbf{statistic}. Parameters are usually unknown;
statistics are what we can compute.

\begin{figure}[H]
\centering
\begin{tikzpicture}
\draw[thick, fill=histfill!15] (0,0) ellipse (3.0cm and 1.7cm);
\node[font=\small] at (0,1.95) {Population: every request today};
\foreach \x/\y in {-2.2/0.3,-1.8/-0.6,-1.5/0.9,-1.2/0.1,-0.9/-1.1,-0.6/0.6,-0.3/-0.4,0/1.2,
  0.2/0.1,0.5/-0.9,0.8/0.7,1.1/-0.2,1.4/1.0,1.6/-0.8,1.9/0.3,2.3/-0.1,-2.4/-0.2,-1.0/1.3,
  0.6/1.25,1.0/-1.25,-0.4/-1.3,2.0/0.9,-1.9/0.85,0.1/-0.7}
  \fill[distblue] (\x,\y) circle (2pt);
\foreach \x/\y in {-1.2/0.1,0.2/0.1,0.8/0.7,1.6/-0.8,-0.3/-0.4,1.9/0.3}
  \draw[trapred, very thick] (\x,\y) circle (4pt);
\draw[thick, fill=quotientorange!15] (7.2,0) circle (1.25cm);
\node[font=\small] at (7.2,1.55) {Sample: what we logged};
\foreach \x/\y in {6.7/0.4,7.3/0.6,7.8/0.1,6.8/-0.4,7.4/-0.2,7.6/-0.7}
  \fill[distblue] (\x,\y) circle (2pt);
\draw[-Stealth, thick] (3.1,0.45) to[bend left=15] node[above, font=\footnotesize] {select} (5.85,0.45);
\draw[-Stealth, thick, densityred] (5.85,-0.45) to[bend left=15] node[below, font=\footnotesize] {infer} (3.1,-0.45);
\end{tikzpicture}
\caption{A sample is a subset of the population. Describing the sample is arithmetic.
Going back along the red arrow, from sample to population, is inference, and it is only as good
as the way the sample was selected.}
\label{fig:data-pop-sample}
\end{figure}

How samples are drawn and why random selection matters is the subject of Chapter 5. For now,
keep the vocabulary straight: in this chapter and the next three we only describe the data in
hand.

\section{Categorical data: counts and bar charts}
\label{sec:data-categorical}

\subsection{Frequency tables}

Fifty requests come back with these status codes: 41 returned 200, 4 returned 304, 3 returned
404, and 2 returned 500. A \textbf{frequency table} lists each category with its count (its
\emph{frequency}). The \textbf{relative frequency} of a category is its count divided by the
total: the fraction of units that fell in that group.

\begin{table}[H]
\centering
\begin{tabular}{lrr}
\toprule
status & frequency & relative frequency\\
\midrule
200 OK            & 41 & $41/50 = 0.82$\\
304 Not Modified  & 4  & $4/50 = 0.08$\\
404 Not Found     & 3  & $3/50 = 0.06$\\
500 Server Error  & 2  & $2/50 = 0.04$\\
\midrule
total             & 50 & $1.00$\\
\bottomrule
\end{tabular}
\caption{Frequency table for 50 status codes. Relative frequencies always sum to 1.}
\label{tab:data-status}
\end{table}

Relative frequencies are what let you compare data sets of different sizes. ``2 errors'' means
nothing until you know whether it was out of 50 requests or 50 million. ``4\% errors'' is
comparable across a quiet minute and a busy hour.

\subsection{Bar charts}

A \textbf{bar chart} draws one bar per category, with height equal to the count or the
relative frequency. The bars are separated by gaps, because the categories are separate
things; there is nothing ``between'' 304 and 404.

\begin{figure}[H]
\centering
\begin{tikzpicture}
\begin{axis}[
  name=left, width=6.3cm, height=5.2cm,
  ybar, bar width=0.55cm, ymin=0, ymax=0.95,
  symbolic x coords={200,304,404,500}, xtick=data,
  ylabel={relative frequency}, xlabel={status code (nominal)},
  title={Nominal: order by size}, title style={font=\small},
  nodes near coords, nodes near coords style={font=\footnotesize},
  every node near coord/.append style={/pgf/number format/fixed, /pgf/number format/precision=2},
  axis lines*=left, enlarge x limits=0.2
]
\addplot[fill=histfill, draw=distblue] coordinates {(200,0.82) (304,0.08) (404,0.06) (500,0.04)};
\end{axis}
\begin{axis}[
  at={(left.east)}, anchor=west, xshift=1.3cm,
  width=6.3cm, height=5.2cm,
  ybar, bar width=0.45cm, ymin=0, ymax=0.42,
  xtick={1,2,3,4,5}, xlabel={star rating (ordinal)},
  title={Ordinal: keep the order}, title style={font=\small},
  nodes near coords, nodes near coords style={font=\footnotesize},
  every node near coord/.append style={/pgf/number format/fixed, /pgf/number format/precision=3},
  axis lines*=left, enlarge x limits=0.22
]
\addplot[fill=structurepurple!40, draw=structurepurple] coordinates {(1,0.075) (2,0.125) (3,0.2) (4,0.35) (5,0.25)};
\end{axis}
\end{tikzpicture}
\caption{Bar charts for categorical data. Left: 50 status codes; with no natural order, sort
bars by height so the big category is obvious. Right: 40 star ratings (counts 3, 5, 8, 14, 10);
the order 1 to 5 carries information, so never re-sort an ordinal axis.}
\label{fig:data-bars}
\end{figure}

\begin{trapbox}
A bar chart whose $y$-axis does not start at zero exaggerates differences. Bar \emph{length} is
what the eye compares. If the axis starts at 0.75, a 0.82 bar is drawn $0.07/0.03\approx 2.3$
times as long as a 0.78 bar, though the two values differ by only 5\%. For bar charts, start
at zero.
\end{trapbox}

\section{Numerical data: seeing every value}
\label{sec:data-small-numerical}

Return to the twenty latencies. With a small data set we can draw every value.

\subsection{Dot plots}

A \textbf{dot plot} puts a number line under the data and stacks one dot per observation
above its value. Nothing is hidden or grouped.

\begin{figure}[H]
\centering
\begin{tikzpicture}
\begin{axis}[
  width=12.5cm, height=3.6cm,
  xmin=0, xmax=150, ymin=0.3, ymax=3.2,
  axis lines*=left, axis y line=none,
  xtick={0,10,...,150}, xlabel={latency (ms)},
  clip=false
]
\addplot[only marks, mark=*, mark size=2.6pt, color=distblue]
  coordinates {(12,1) (14,1) (15,1) (15,2) (17,1) (18,1) (18,2) (19,1) (21,1) (22,1) (23,1) (23,2) (24,1) (26,1) (28,1) (31,1) (35,1) (42,1) (58,1)};
\addplot[only marks, mark=*, mark size=2.6pt, color=trapred] coordinates {(140,1)};
\draw[decorate, decoration={brace, amplitude=5pt}] (axis cs:11,2.5) -- (axis cs:29,2.5)
  node[midway, above=5pt, font=\footnotesize] {15 of 20 values};
\draw[<->, gray] (axis cs:61,1.7) -- (axis cs:137,1.7) node[midway, above, font=\footnotesize] {empty gap of 82 ms};
\end{axis}
\end{tikzpicture}
\caption{Dot plot of the 20 latencies. Most values crowd between 12 and 28 ms, a few trail
off to the right, and one sits far away at 140 ms. This is the whole data set; every dot is a
request.}
\label{fig:data-dotplot}
\end{figure}

Three features are visible at once that were invisible in the list: where the bulk of the data
sits (around 20 ms), that the values trail off further to the right than to the left, and that
140 is separated from everything else by a wide gap.

\subsection{Stem-and-leaf plots}

A \textbf{stem-and-leaf plot} is a dot plot made of digits. Split each number into a
\emph{stem} (all digits but the last) and a \emph{leaf} (the last digit). For 23, the stem is 2
and the leaf is 3. Write each stem once, then list the leaves of all values with that stem.

\begin{figure}[H]
\centering
\begin{tikzpicture}
\node[draw=gray!60, fill=codebg, inner sep=8pt] (s) {%
\ttfamily
\begin{tabular}{r|l}
1 & 2 4 5 5 7 8 8 9\\
2 & 1 2 3 3 4 6 8\\
3 & 1 5\\
4 & 2\\
5 & 8\\
6 & \\
7 & \\
8 & \\
9 & \\
10 & \\
11 & \\
12 & \\
13 & \\
14 & 0
\end{tabular}};
\node[right=0.6cm of s, align=left, font=\small, text width=5.6cm] {Key: \texttt{2 | 3} means 23 ms.\\[4pt]
Turn the page sideways and the leaves form a histogram with bins of width 10. Unlike a
histogram, you can read back every original value.\\[4pt]
The run of empty stems from 6 to 13 is the gap before the outlier.};
\end{tikzpicture}
\caption{Stem-and-leaf plot of the 20 latencies. It keeps every digit and shows the shape at
the same time, which makes it a good hand tool for small data sets.}
\label{fig:data-stemleaf}
\end{figure}

Stem-and-leaf plots stop being practical beyond about a hundred values. Their real use is
reading small data by hand, and they teach the move that the histogram makes: group nearby
values together and count.

\section{Grouping numerical data: histograms}
\label{sec:data-histograms}

\subsection{From frequency table to histogram}

For a numerical variable, a frequency table needs \textbf{bins}: intervals that split the
number line into adjacent, non-overlapping pieces. Use bins of width 10 ms, closed on the left
and open on the right, so $[10,20)$ contains 10 but not 20. Count the latencies in each bin.

\begin{table}[H]
\centering
\begin{tabular}{lcrr}
\toprule
bin (ms) & values & frequency & relative frequency\\
\midrule
$[10,20)$ & 12 14 15 15 17 18 18 19 & 8 & 0.40\\
$[20,30)$ & 21 22 23 23 24 26 28 & 7 & 0.35\\
$[30,40)$ & 31 35 & 2 & 0.10\\
$[40,50)$ & 42 & 1 & 0.05\\
$[50,60)$ & 58 & 1 & 0.05\\
$[60,140)$ & --- & 0 & 0.00\\
$[140,150)$ & 140 & 1 & 0.05\\
\midrule
total & & 20 & 1.00\\
\bottomrule
\end{tabular}
\caption{Binned frequency table for the 20 latencies. The eight bins from 60 to 140 are
empty and have been collapsed into one row.}
\label{tab:data-latency-bins}
\end{table}

A \textbf{histogram} draws this table. Each bin becomes a bar sitting on the number line, as
wide as the bin, as tall as the count. Adjacent bars touch, because the bins tile the number
line with no space between them.

\begin{figure}[H]
\centering
\begin{tikzpicture}
\begin{axis}[
  width=12.5cm, height=5cm,
  ymin=0, ymax=9.5, xmin=0, xmax=150,
  xtick={0,10,...,150}, ytick={0,2,...,8},
  xlabel={latency (ms)}, ylabel={count},
  axis lines*=left
]
\addplot[ybar interval, fill=histfill, draw=distblue]
  coordinates {(10,8) (20,7) (30,2) (40,1) (50,1) (60,0) (70,0) (80,0) (90,0) (100,0) (110,0) (120,0) (130,0) (140,1) (150,0)};
\end{axis}
\end{tikzpicture}
\caption{Histogram of the 20 latencies with bins of width 10. The bars touch because latency
is continuous: every value between 10 and 150 belongs to exactly one bin.}
\label{fig:data-hist20}
\end{figure}

\begin{intuitionbox}
Bar chart or histogram? Ask what the horizontal axis is. If it is a list of labels, you can
shuffle them, and gaps between bars say ``these are separate things'': bar chart. If it is a
number line, position matters and there are no gaps in the number line, so there are none
between bars: histogram.
\end{intuitionbox}

The same tables in Python, for the status codes and the 20 latencies:

\begin{lstlisting}[language=Python]
from collections import Counter
import numpy as np
codes = [200]*41 + [304]*4 + [404]*3 + [500]*2
print(Counter(codes).most_common())      # categorical: counts
lat = np.array([12,14,15,15,17,18,18,19,21,22,23,23,24,26,28,31,35,42,58,140])
counts, edges = np.histogram(lat, bins=np.arange(10, 151, 10))
print(counts)                            # [8 7 2 1 1 0 0 0 0 0 0 0 0 1]
print(counts / len(lat))                 # relative frequencies
\end{lstlisting}

\subsection{Discrete numerical data: coin flips}

Discrete numerical data sits between bar charts and histograms: the axis is a number line,
but only whole numbers can occur. Flip a coin 20 times:
\[
\texttt{H T H H T T H T H H H T T H T H H T H T}
\]
The outcome of one flip is categorical (heads or tails); its frequency table is 11 heads,
9 tails, relative frequencies 0.55 and 0.45. Now change the unit. Let one unit be \emph{a
batch of 10 flips}, and the variable be \emph{the number of heads in the batch}. That variable
is discrete numerical: it takes the values $0,1,\dots,10$. Repeat for 100 batches and count.

\begin{figure}[H]
\centering
\begin{tikzpicture}
\begin{axis}[
  name=coinL, width=4.4cm, height=5cm,
  ybar, bar width=0.6cm, ymin=0, ymax=14,
  symbolic x coords={H,T}, xtick=data, enlarge x limits=0.6,
  xlabel={one flip}, ylabel={count},
  title={20 single flips}, title style={font=\small},
  nodes near coords, nodes near coords style={font=\footnotesize},
  axis lines*=left
]
\addplot[fill=quotientorange!50, draw=quotientorange] coordinates {(H,11) (T,9)};
\end{axis}
\begin{axis}[
  at={(coinL.east)}, anchor=west, xshift=1.2cm,
  width=8cm, height=5cm,
  ybar, bar width=0.32cm, ymin=0, ymax=29,
  xmin=-0.6, xmax=10.6, xtick={0,1,...,10},
  xlabel={heads in a batch of 10}, ylabel={batches},
  title={100 batches of 10 flips}, title style={font=\small},
  nodes near coords, nodes near coords style={font=\footnotesize},
  axis lines*=left
]
\addplot[fill=histfill, draw=distblue] coordinates {(0,0) (1,0) (2,4) (3,13) (4,19) (5,24) (6,20) (7,14) (8,3) (9,3) (10,0)};
\end{axis}
\end{tikzpicture}
\caption{Left: the outcome of one flip is categorical. Right: the number of heads in 10 flips
is discrete numerical. Bars sit on the integers, thin enough that you cannot mistake 4.5 for a
possible value. The pile centers near 5 and falls off roughly symmetrically on both sides.}
\label{fig:data-coins}
\end{figure}

Exactly 5 heads happened in $24$ of the $100$ batches. Probability theory predicts
$\binom{10}{5}/2^{10} = 252/1024 \approx 0.246$ (see \emph{Probability From the Ground Up},
Ch.~7 on the binomial). The observed 0.24 is close but not equal. That gap between a
relative frequency and the probability behind it is the gap that inference is about.

\subsection{Larger data: 200 latencies}

Dot plots and stem-and-leaf plots break down with many values. Here are 200 latencies from
the same kind of endpoint, ranging from 13 to 110 ms. Grouped into bins of width 25:

\begin{center}
\begin{tabular}{lccccc}
\toprule
bin (ms) & $[0,25)$ & $[25,50)$ & $[50,75)$ & $[75,100)$ & $[100,125)$\\
\midrule
count & 29 & 120 & 43 & 7 & 1\\
relative frequency & 0.145 & 0.600 & 0.215 & 0.035 & 0.005\\
\bottomrule
\end{tabular}
\end{center}

These counts are correct. Whether they are \emph{useful} depends on the bins, and that is
the next problem.

\section{Bin width changes the picture}
\label{sec:data-binwidth}

There is no single histogram of a data set. Every choice of bin width gives a different
picture. Figure~\ref{fig:data-binwidths} shows the same 200 latencies three ways.

\begin{figure}[H]
\centering
\begin{tikzpicture}
\begin{axis}[
  name=b1, width=12.5cm, height=3.9cm,
  ymin=0, ymax=10, xmin=0, xmax=125, xtick={0,25,...,125}, xticklabels={},
  ylabel={count}, axis lines*=left,
  title={width 1 ms: too narrow, mostly noise}, title style={font=\small}
]
\addplot[ybar interval, fill=histfill, draw=distblue, line width=0.2pt]
  coordinates {(13,1) (14,0) (15,1) (16,3) (17,2) (18,2) (19,1) (20,3) (21,2) (22,5) (23,6) (24,3) (25,3) (26,4) (27,4) (28,6) (29,2) (30,7) (31,8) (32,6) (33,6) (34,3) (35,8) (36,4) (37,8) (38,6) (39,8) (40,8) (41,6) (42,7) (43,2) (44,3) (45,2) (46,2) (47,5) (48,1) (49,1) (50,2) (51,2) (52,2) (53,1) (54,6) (55,3) (56,3) (57,2) (58,1) (59,2) (60,2) (61,2) (62,3) (63,0) (64,2) (65,0) (66,1) (67,1) (68,2) (69,0) (70,3) (71,0) (72,0) (73,2) (74,1) (75,1) (76,1) (77,1) (78,0) (79,2) (80,0) (81,0) (82,0) (83,0) (84,1) (85,0) (86,0) (87,0) (88,0) (89,0) (90,0) (91,0) (92,0) (93,0) (94,0) (95,0) (96,0) (97,0) (98,1) (99,0) (100,0) (101,0) (102,0) (103,0) (104,0) (105,0) (106,0) (107,0) (108,0) (109,0) (110,1) (111,0)};
\end{axis}
\begin{axis}[
  name=b5, at={(b1.south)}, anchor=north, yshift=-1.15cm,
  width=12.5cm, height=3.9cm,
  ymin=0, ymax=38, xmin=0, xmax=125, xtick={0,25,...,125}, xticklabels={},
  ylabel={count}, axis lines*=left,
  title={width 5 ms: shape is clear}, title style={font=\small}
]
\addplot[ybar interval, fill=histfill, draw=distblue]
  coordinates {(10,1) (15,9) (20,19) (25,19) (30,30) (35,34) (40,26) (45,11) (50,13) (55,11) (60,9) (65,4) (70,6) (75,5) (80,1) (85,0) (90,0) (95,1) (100,0) (105,0) (110,1) (115,0)};
\end{axis}
\begin{axis}[
  name=b25, at={(b5.south)}, anchor=north, yshift=-1.15cm,
  width=12.5cm, height=3.9cm,
  ymin=0, ymax=135, xmin=0, xmax=125, xtick={0,25,...,125},
  xlabel={latency (ms)}, ylabel={count}, axis lines*=left,
  title={width 25 ms: too wide, detail lost}, title style={font=\small}
]
\addplot[ybar interval, fill=histfill, draw=distblue]
  coordinates {(0,29) (25,120) (50,43) (75,7) (100,1) (125,0)};
\end{axis}
\end{tikzpicture}
\caption{The same 200 latencies with three bin widths. Width 1 shows the random jitter of
individual counts (note the meaningless spikes and holes). Width 25 squeezes the whole data
set into five bars and hides the long right tail. Width 5 shows a single peak near 35 ms and a
tail stretching to the right.}
\label{fig:data-binwidths}
\end{figure}

What goes wrong at each extreme:

\begin{itemize}
  \item \textbf{Too narrow.} Each bin holds only a few values, so its count is dominated by
  luck. The bar at 31 ms has height 8, the bar at 34 ms has height 3. Nobody believes requests
  ``prefer'' 31 ms over 34 ms; another sample would put the spikes elsewhere.
  \item \textbf{Too wide.} Each bin averages over a big stretch of the number line. The
  $[25,50)$ bar says 60\% of requests fall there but cannot tell you whether they cluster near
  30 or near 45. The tail beyond 75 ms, which is exactly what an on-call engineer cares about,
  shrinks to two stubby bars.
\end{itemize}

\begin{whybox}
Bin width trades two errors against each other. A narrow bin is \emph{faithful to its location}
but holds few points, so its height is noisy. A wide bin holds many points, so its height is
stable, but it blurs everything inside it into one flat bar. With more data, each bin can be
narrower before noise takes over. That is why rules of thumb for the number of bins grow with
$n$: Sturges' rule uses $1+\log_2 n$ bins, which for $n=200$ gives about $8.6$, and the
Freedman--Diaconis rule sets the width to $2\cdot\IQR\cdot n^{-1/3}$ (the $\IQR$, a measure of
spread, is defined in Chapter 2). For the 200 latencies, which run from 13 to 110 ms, Sturges
suggests a width of about $97/8.6\approx 11$ ms and Freedman--Diaconis about 7 ms. Treat them
as starting points and always try a few widths.
\end{whybox}

\begin{trapbox}
Changing the bin width or shifting where the bins start can make a bump appear or vanish.
If a histogram is evidence for some claim (``there are two kinds of users''), check that the
feature survives at two or three other bin widths before you believe it.
\end{trapbox}

\section{Density histograms: area is proportion}
\label{sec:data-density}

\subsection{The problem with unequal bins}

Sometimes bins must have different widths. Dashboards often use wide bins in the tail, where
data is sparse, and narrow bins in the middle. Take the 200 latencies with edges
$0, 20, 30, 40, 50, 60, 80, 120$. The counts are $10, 38, 64, 37, 24, 24, 3$. Plot those counts
as bar heights and you get the left panel of Figure~\ref{fig:data-unequal}.

\begin{figure}[H]
\centering
\begin{tikzpicture}
\begin{axis}[
  name=uL, width=6.4cm, height=5.2cm,
  ymin=0, ymax=70, xmin=0, xmax=120, xtick={0,20,...,120},
  xlabel={latency (ms)}, ylabel={count},
  title={Height = count (misleading)}, title style={font=\small},
  axis lines*=left
]
\addplot[ybar interval, fill=trapred!35, draw=trapred]
  coordinates {(0,10) (20,38) (30,64) (40,37) (50,24) (60,24) (80,3) (120,0)};
\end{axis}
\begin{axis}[
  at={(uL.east)}, anchor=west, xshift=1.5cm,
  width=6.4cm, height=5.2cm,
  ymin=0, ymax=0.035, xmin=0, xmax=120, xtick={0,20,...,120},
  xlabel={latency (ms)}, ylabel={density (per ms)},
  title={Height = density (honest)}, title style={font=\small},
  yticklabel style={/pgf/number format/fixed, /pgf/number format/precision=2},
  scaled y ticks=false, axis lines*=left
]
\addplot[ybar interval, fill=histfill, draw=distblue]
  coordinates {(0,0.0025) (20,0.019) (30,0.032) (40,0.0185) (50,0.012) (60,0.006) (80,0.000375) (120,0)};
\end{axis}
\end{tikzpicture}
\caption{Same data, same unequal bins. Left: the $[60,80)$ bar is as tall as the $[50,60)$
bar because both hold 24 requests, but its 24 are spread over twice the width, so the data is
half as dense there. The wide $[80,120)$ bar has a big visible area for just 3 requests. Right:
dividing by bin width fixes both distortions.}
\label{fig:data-unequal}
\end{figure}

The eye reads a bar's \emph{area}, not its height. A bar twice as wide looks like twice as
much data. The fix is to choose heights so that area tells the truth.

\begin{definition}[Density histogram]
For a data set of $n$ values and a bin of width $w$ containing $k$ values, the bar height is
\[
  \text{density} = \frac{k}{n\,w} = \frac{\text{relative frequency}}{\text{bin width}}.
\]
The units are ``fraction of the data per unit of the variable,'' here per millisecond.
\end{definition}

For the bin $[60,80)$: $k=24$, $n=200$, $w=20$, so density $= 24/(200\cdot 20) = 0.006$ per ms.
For $[50,60)$: $24/(200\cdot 10) = 0.012$. Same count, half the density, as it should be.

\begin{whybox}
The area of a bar is height times width:
\[
  \frac{k}{n\,w}\cdot w = \frac{k}{n},
\]
which is exactly the fraction of the data in that bin. Since every value lands in exactly one
bin, the bar areas add up to $\sum k/n = n/n = 1$. Check it on our bins:
\[
  0.05 + 0.19 + 0.32 + 0.185 + 0.12 + 0.12 + 0.015 = 1.
\]
So on a density histogram, ``what fraction of requests took between $a$ and $b$ ms?'' is
answered by ``the area of the bars between $a$ and $b$,'' whatever the bin widths.
\end{whybox}

\begin{figure}[H]
\centering
\begin{tikzpicture}
\begin{axis}[
  width=12.5cm, height=5.4cm,
  ymin=0, ymax=0.04, xmin=0, xmax=125, xtick={0,10,...,120},
  xlabel={latency (ms)}, ylabel={density (per ms)},
  yticklabel style={/pgf/number format/fixed, /pgf/number format/precision=2},
  scaled y ticks=false, axis lines*=left, clip=false
]
\addplot[ybar interval, fill=histfill!50, draw=distblue]
  coordinates {(10,0.001) (15,0.009) (20,0.019) (25,0.019) (30,0.030) (35,0.034) (40,0.026) (45,0.011) (50,0.013) (55,0.011) (60,0.009) (65,0.004) (70,0.006) (75,0.005) (80,0.001) (85,0) (90,0) (95,0.001) (100,0) (105,0) (110,0.001) (115,0)};
\addplot[ybar interval, fill=quotientorange!70, draw=quotientorange!80!black]
  coordinates {(30,0.030) (35,0.034) (40,0)};
\node[font=\small, align=left, anchor=west] at (axis cs:58,0.031)
  {shaded area $= 0.030\cdot 5 + 0.034\cdot 5 = 0.32$\\ so 32\% of requests took 30--40 ms};
\draw[-Stealth, quotientorange!80!black] (axis cs:57,0.031) -- (axis cs:41,0.029);
\end{axis}
\end{tikzpicture}
\caption{Density histogram of the 200 latencies, bins of width 5. The orange area equals
the fraction of requests between 30 and 40 ms. Counting directly confirms it: 64 of the 200
latencies lie in $[30,40)$, and $64/200 = 0.32$.}
\label{fig:data-density-area}
\end{figure}

\begin{connectionbox}
Area-equals-proportion is the idea that turns a histogram into a probability density. As the
data grows and the bins shrink, the tops of the bars trace out a smooth curve $f(x)$, and
``fraction of data between $a$ and $b$'' becomes $\int_a^b f(x)\,\dd x$. That is precisely the
definition of a probability density function (\emph{Probability From the Ground Up}, Ch.~10).
When an ML paper writes $x \sim p$, read it as ``$x$ is drawn from $p$'': the $p$ is this curve.
\end{connectionbox}

\begin{lstlisting}[language=Python]
import numpy as np
lat = np.loadtxt("latencies.txt")          # the 200 values
edges = [0, 20, 30, 40, 50, 60, 80, 120]
counts, _ = np.histogram(lat, bins=edges)   # [10 38 64 37 24 24 3]
dens, _ = np.histogram(lat, bins=edges, density=True)
print(dens)                                 # 0.0025 0.019 0.032 ...
print((dens * np.diff(edges)).sum())        # 1.0: areas sum to one
\end{lstlisting}

\begin{recoverybox}
A density histogram has a bar over $[40,50)$ of height $0.0185$ per ms, with $n = 200$. How
many values are in that bin? \emph{Answer:} area $= 0.0185\times 10 = 0.185$, and
$0.185 \times 200 = 37$.
\end{recoverybox}

\section{Shapes of distributions}
\label{sec:data-shapes}

Once the bins are sensible, the first thing to read off a histogram is its overall shape.
A handful of shapes cover most of what you will meet.

\begin{figure}[H]
\centering
\begin{tikzpicture}
\pgfplotsset{shp/.style={width=4.6cm, height=3.5cm, xmin=0, xmax=10, ymin=0, ymax=1.15,
  axis lines*=left, xtick=\empty, ytick=\empty, title style={font=\small}, domain=0:10, samples=120}}
\begin{axis}[shp, name=s1, title={symmetric}]
\addplot[fill=histfill!50, draw=distblue, thick] {exp(-(x-5)^2/2)} \closedcycle;
\end{axis}
\begin{axis}[shp, name=s2, at={(s1.east)}, anchor=west, xshift=0.5cm, title={skewed right}]
\addplot[fill=histfill!50, draw=distblue, thick] {x*exp(1-x/1.5)/1.5} \closedcycle;
\end{axis}
\begin{axis}[shp, name=s3, at={(s2.east)}, anchor=west, xshift=0.5cm, title={skewed left}]
\addplot[fill=histfill!50, draw=distblue, thick] {(10-x)*exp(1-(10-x)/1.5)/1.5} \closedcycle;
\end{axis}
\begin{axis}[shp, name=s4, at={(s1.south)}, anchor=north, yshift=-1.25cm, xshift=3.8cm, title={uniform}, samples=200]
\addplot[fill=histfill!50, draw=distblue, thick] {(x>1)*(x<9)*0.8} \closedcycle;
\end{axis}
\begin{axis}[shp, at={(s4.east)}, anchor=west, xshift=0.5cm, title={bimodal}]
\addplot[fill=histfill!50, draw=distblue, thick] {exp(-(x-3)^2/0.9) + 0.7*exp(-(x-7.2)^2/0.9)} \closedcycle;
\end{axis}
\end{tikzpicture}
\caption{Five basic shapes, drawn as smoothed histogram outlines. Skew is named for the
direction of the long tail, not the peak: ``skewed right'' has its peak on the left and a tail
reaching right.}
\label{fig:data-shapes}
\end{figure}

\begin{itemize}
  \item \textbf{Symmetric}: the left half mirrors the right half. Heights of adults of one sex,
  measurement errors, and the coin-flip counts in Figure~\ref{fig:data-coins} look like this.
  \item \textbf{Skewed right} (positively skewed): a long tail toward large values. Latencies,
  incomes, file sizes, house prices. Anything with a hard floor (latency cannot go below
  zero, or below the network round trip) but no hard ceiling tends to skew right.
  \item \textbf{Skewed left} (negatively skewed): a long tail toward small values. Exam scores
  on an easy test, age at death in a rich country. A hard ceiling with no floor nearby produces
  left skew.
  \item \textbf{Uniform}: roughly flat, every region equally common. The last digit of a
  timestamp in ms, or the face of a fair die over many rolls.
  \item \textbf{Bimodal}: two peaks. Almost always a sign that two different groups have been
  mixed into one column.
\end{itemize}

A peak is called a \textbf{mode}; ``unimodal'' means one peak and ``bimodal'' means two.

\begin{trapbox}
``Skewed right'' confuses people because the bulk of the data is on the \emph{left}. Name the
skew after the tail: if the long thin end points toward large values, the distribution is
skewed right.
\end{trapbox}

\subsection{A left-skewed example: exam scores}

Eighty students take an exam scored out of 100. Their scores, in bins of 10, are:
$[30,40)$: 1, $[40,50)$: 2, $[50,60)$: 2, $[60,70)$: 5, $[70,80)$: 17, $[80,90)$: 31,
$[90,100]$: 22. (The top bin includes 100.)

\begin{figure}[H]
\centering
\begin{tikzpicture}
\begin{axis}[
  width=11cm, height=5cm,
  ymin=0, ymax=36, xmin=30, xmax=100, xtick={30,40,...,100},
  xlabel={exam score}, ylabel={students},
  axis lines*=left, clip=false
]
\addplot[ybar interval, fill=targetgreen!35, draw=targetgreen!70!black]
  coordinates {(30,1) (40,2) (50,2) (60,5) (70,17) (80,31) (90,22) (100,0)};
\draw[densityred, very thick] (axis cs:100,0) -- (axis cs:100,33);
\node[densityred, font=\footnotesize, anchor=south] at (axis cs:100,33) {ceiling};
\draw[-Stealth, gray, thick] (axis cs:62,9) -- (axis cs:36,4);
\node[font=\footnotesize, anchor=south] at (axis cs:52,12) {long tail to the left};
\end{axis}
\end{tikzpicture}
\caption{Exam scores of 80 students. The test was easy for most: the peak is in the 80s and
nobody can score above 100. A few struggling students form a long tail toward low scores,
making the distribution skewed left.}
\label{fig:data-exam}
\end{figure}

Here $53$ of $80$ students ($66\%$) scored 80 or above, while the lowest bin reaches down into
the 30s. Chapter 2 shows that a tail like this pulls the average below the middle value: in a
left-skewed distribution the mean sits below the median.

\subsection{A bimodal example: cache hits and misses}

A service answers from a cache when it can (fast) and from a database when it cannot (slow).
Here are 150 latencies, bins of width 5 ms.

\begin{figure}[H]
\centering
\begin{tikzpicture}
\begin{axis}[
  width=12.5cm, height=5.2cm,
  ymin=0, ymax=58, xmin=0, xmax=110, xtick={0,10,...,110},
  xlabel={latency (ms)}, ylabel={count},
  axis lines*=left, clip=false
]
\addplot[ybar interval, fill=histfill, draw=distblue]
  coordinates {(10,9) (15,49) (20,41) (25,5) (30,1) (35,0) (40,0) (45,1) (50,0) (55,1) (60,4) (65,8) (70,4) (75,11) (80,5) (85,3) (90,3) (95,3) (100,2) (105,0)};
\node[font=\footnotesize, anchor=west] at (axis cs:27,46) {cache hits (fast)};
\node[font=\footnotesize, anchor=south] at (axis cs:82,14) {cache misses (slow)};
\draw[densityred, dashed, thick] (axis cs:37,0) -- (axis cs:37,32);
\node[densityred, font=\footnotesize, anchor=west, align=left] at (axis cs:38.5,30) {mean $\approx 37$ ms:\\no request is\\near this value};
\end{axis}
\end{tikzpicture}
\caption{A bimodal latency distribution. Two kinds of request share one column. The overall
mean lands in the valley between the peaks, a value that describes almost none of the actual
requests.}
\label{fig:data-bimodal}
\end{figure}

Of the 150 requests, only 2 took between 30 and 55 ms, and none between 35 and 45. Yet the
average of all 150 is about 37 ms. The arithmetic shows why: 105 hits (70\%) average about
20 ms and 45 misses (30\%) average about 78 ms, so the overall mean is roughly
$0.7\cdot 20 + 0.3\cdot 78 \approx 37$, a weighted point between the two peaks. A single
summary number cannot describe two groups. The right move is to split the
data (by a cache-hit flag, if the log has one) and describe each group separately. A
histogram is how you notice that you need to.

\section{Outliers}
\label{sec:data-outliers}

An \textbf{outlier} is a value that sits far from the bulk of the data, separated from it by a
gap. The 140 ms request in Figure~\ref{fig:data-dotplot} is one: the other 19 values are all
under 60, and nothing lies between 58 and 140.

There is no universal cutoff for ``far''; Chapter 3 gives a standard rule based on quartiles.
What matters more is what you \emph{do} with an outlier. It can be any of these:

\begin{itemize}
  \item \textbf{A recording error.} A latency of $-3$ ms, a height of 17 m, a score of 850 on a
  100-point test. Fix it if you can, drop it if you cannot, and write down that you did.
  \item \textbf{A different process.} The 140 ms request returned 503 and retried three times
  (Table~\ref{tab:data-log}). It is a real request, but it is a failure, not a slow success.
  \item \textbf{The thing you care about.} For a latency service-level objective (a promise
  such as ``99\% of requests finish within 100 ms''), the slow
  requests \emph{are} the story. Deleting them would make the dashboard look better and the
  product worse.
\end{itemize}

Outliers can also come in clumps. Clients usually give up after a fixed timeout, and the log
records the cap rather than the true time.

\begin{figure}[H]
\centering
\begin{tikzpicture}
\begin{axis}[
  width=12.5cm, height=4.8cm,
  ymin=0, ymax=135, xmin=0, xmax=1050, xtick={0,100,...,1000},
  xlabel={latency (ms)}, ylabel={count},
  axis lines*=left, clip=false
]
\addplot[ybar interval, fill=histfill, draw=distblue]
  coordinates {(0,29) (25,120) (50,43) (75,7) (100,1) (125,0)};
\addplot[ybar interval, fill=trapred!60, draw=trapred]
  coordinates {(1000,6) (1025,0)};
\node[font=\footnotesize, anchor=south east, align=right, text=trapred] at (axis cs:1010,14)
  {6 requests hit the\\1000 ms client timeout};
\node[font=\footnotesize, anchor=west, align=left] at (axis cs:140,100)
  {200 completed requests,\\all under 125 ms};
\end{axis}
\end{tikzpicture}
\caption{The 200 latencies plus 6 timeouts recorded as exactly 1000 ms. The red bar is not
part of the tail; it is a pile-up at a hard cap, and the true times of those requests are
unknown (they are at least 1000 ms). Averaging it in would treat ``gave up'' as a measurement.}
\label{fig:data-timeouts}
\end{figure}

\begin{trapbox}
Never delete outliers just because they are inconvenient. Investigate first. Whatever you
decide, report results both with and without them if the conclusion changes.
\end{trapbox}

\section{The distribution is the object of study}
\label{sec:data-distribution}

Every chart in this chapter answers the same question: \emph{which values occur, and how
often?} That pairing of values with their frequencies is the distribution.

\begin{definition}[Distribution of a variable]
The \textbf{distribution} of a variable in a data set is the list of values it takes together
with how often it takes each one (as counts, relative frequencies, or densities). For a
population, the distribution describes the whole population; for a sample, it is called the
\textbf{empirical distribution}.
\end{definition}

The empirical distribution is a rough copy of the population distribution, and the copy gets
sharper as the sample grows. Figure~\ref{fig:data-growth} shows density histograms of latency
samples of three sizes, all drawn from the same underlying process, with the population density
drawn on top.

\begin{figure}[H]
\centering
\begin{tikzpicture}
\pgfplotsset{grw/.style={width=4.55cm, height=4.2cm, xmin=0, xmax=150, ymin=0, ymax=0.03,
  xtick={0,50,100}, ytick=\empty, axis lines*=left, title style={font=\small},
  xlabel={ms}, xlabel style={font=\footnotesize}, ticklabel style={font=\footnotesize}}}
\begin{axis}[grw, name=g1, title={$n=50$}]
\addplot[ybar interval, fill=histfill!60, draw=distblue]
  coordinates {(0,0) (10,0.006) (20,0.022) (30,0.024) (40,0.020) (50,0.012) (60,0.010) (70,0.002) (80,0.002) (90,0) (100,0) (110,0) (120,0) (130,0.002) (140,0) (150,0)};
\addplot[densityred, thick, domain=2:150, samples=120] {0.88654/x*exp(-(ln(x)-3.688879)^2/0.405)};
\end{axis}
\begin{axis}[grw, name=g2, at={(g1.east)}, anchor=west, xshift=0.5cm, title={$n=500$}]
\addplot[ybar interval, fill=histfill!60, draw=distblue]
  coordinates {(0,0) (10,0.0048) (20,0.0202) (30,0.0234) (40,0.0178) (50,0.0152) (60,0.0068) (70,0.0046) (80,0.0036) (90,0.0010) (100,0.0010) (110,0.0006) (120,0) (130,0) (140,0.0006) (150,0)};
\addplot[densityred, thick, domain=2:150, samples=120] {0.88654/x*exp(-(ln(x)-3.688879)^2/0.405)};
\end{axis}
\begin{axis}[grw, at={(g2.east)}, anchor=west, xshift=0.5cm, title={$n=5000$}]
\addplot[ybar interval, fill=histfill!60, draw=distblue]
  coordinates {(0,0.0002) (10,0.0058) (20,0.0190) (30,0.0251) (40,0.0188) (50,0.0127) (60,0.0073) (70,0.0047) (80,0.0024) (90,0.0015) (100,0.0008) (110,0.0008) (120,0.0003) (130,0.0001) (140,0.0002) (150,0)};
\addplot[densityred, thick, domain=2:150, samples=120] {0.88654/x*exp(-(ln(x)-3.688879)^2/0.405)};
\end{axis}
\end{tikzpicture}
\caption{Density histograms of three samples from one latency process, bin width 10 ms, all on
the same scale. The red curve is the population density. At $n=50$ the bars are lumpy; at
$n=5000$ they hug the curve. The population distribution is the fixed target; each sample is a
noisy picture of it.}
\label{fig:data-growth}
\end{figure}

\begin{happeningbox}
This figure is a preview of the whole book. The red curve is what we want to know and never
see directly. The bars are what we have. Describing the bars is Chapters 1--4. Saying how far
the bars might be from the curve, for a given $n$, is inference: sampling distributions,
standard errors, and confidence intervals (Chapters 6--8). Using a fitted curve to say what
comes next is prediction (Chapters 12--15).
\end{happeningbox}

\subsection{How to read any distribution}

When you meet a new histogram, read it in a fixed order. Each step points to the chapter that
develops it.

\begin{figure}[H]
\centering
\begin{tikzpicture}[
  st/.style={draw, rounded corners, thick, fill=histfill!20, align=center, font=\small, minimum height=1.3cm, text width=2.35cm},
  arr/.style={-Stealth, thick}
]
\node[st] (a) at (0,0) {\textbf{Shape}\\one peak or two?\\skewed?};
\node[st] (b) at (3.1,0) {\textbf{Center}\\where is a typical value?};
\node[st] (c) at (6.2,0) {\textbf{Spread}\\how far do values range?};
\node[st] (d) at (9.3,0) {\textbf{Unusual}\\gaps, outliers, pile-ups?};
\draw[arr] (a) -- (b);
\draw[arr] (b) -- (c);
\draw[arr] (c) -- (d);
\node[font=\footnotesize, below=0.15cm of a] {this chapter};
\node[font=\footnotesize, below=0.15cm of b] {Chapter 2};
\node[font=\footnotesize, below=0.15cm of c] {Chapters 2--3};
\node[font=\footnotesize, below=0.15cm of d] {Chapters 1, 3};
\end{tikzpicture}
\caption{A fixed reading order for any distribution. Shape comes first because it decides
which measures of center and spread are trustworthy: the mean works for the symmetric coin
counts and fails for the bimodal cache latencies.}
\label{fig:data-reading-order}
\end{figure}

Apply it to the 200 latencies (Figure~\ref{fig:data-binwidths}, middle panel). \emph{Shape}:
one peak, skewed right. \emph{Center}: the peak bin is $[35,40)$ and half the values are at or
below 38 ms. \emph{Spread}: most values lie between 15 and 80 ms. \emph{Unusual}: isolated
values at 98 and 110 ms worth a look. Those four sentences already tell an on-call engineer
what to expect from this endpoint.

\begin{historybox}
Karl Pearson, a British mathematician building the new field of mathematical statistics,
introduced the word \emph{histogram} in his lectures in 1891. He needed a name for the
bar-picture he used to compare measured data (skull sizes, crab shell widths, death ages)
with the smooth curves his theory predicted, and the name stuck. Eighty years later, John
Tukey, a Princeton and Bell Labs statistician, complained that statisticians rushed to fit
models before looking at their data. In the 1970s he invented quick pencil-and-paper displays
for looking first, including the stem-and-leaf plot and the box plot (Chapter 3), and collected
them in his 1977 book \emph{Exploratory Data Analysis}. His point became standard practice:
plot the distribution before you compute anything.
\end{historybox}


\section{Summary}
\label{sec:data-summary}

\begin{itemize}
  \item Statistics describes data in hand, infers facts about a population, and predicts new
  values. Description comes first.
  \item A data set is a table: rows are observational units, columns are variables. Decide
  what one row is before analyzing.
  \item Variables are categorical (nominal or ordinal) or numerical (discrete or continuous).
  Digits do not make a variable numerical; ask whether arithmetic makes sense.
  \item A sample is the part of the population you observed. Parameters describe populations;
  statistics describe samples.
  \item Frequency tables count; relative frequencies (count$/n$) sum to 1 and compare across
  data sets of different sizes.
  \item Categorical data: bar charts with gaps, starting at zero; keep ordinal order.
  Numerical data: dot plots and stem-and-leaf for small $n$, histograms with touching bars for
  larger $n$.
  \item Bin width trades noise (too narrow) against blur (too wide). Try several.
  \item Density histogram: height $= k/(nw)$, so bar area is the fraction of data in the bin and
  total area is 1. Required for unequal bins; the bridge to probability densities.
  \item Shapes: symmetric, skewed right (long right tail), skewed left, uniform, bimodal.
  Bimodal usually means two groups mixed together.
  \item Outliers may be errors, a different process, or the main story. Investigate before
  removing.
  \item The distribution (which values, how often) is the object of study. Read it in order:
  shape, center, spread, unusual features.
\end{itemize}

Chapter 2 turns the second and third items of that reading order into numbers: the mean and
median for center, and the standard deviation for spread, including why it divides by $n-1$.

\section{Exercises}
\label{sec:data-exercises}

\begin{enumerate}
  \item For each variable, give its type (nominal, ordinal, discrete, continuous): (a) CPU
  temperature; (b) number of open pull requests; (c) country of a user; (d) severity level
  P0/P1/P2/P3; (e) HTTP method (GET, POST, \dots).
  \item A survey of 250 engineers asks their primary language: Python 110, Java 60, Go 45,
  C\texttt{++} 35. Compute the relative frequencies and say which chart you would draw and how
  you would order its bars.
  \item Make a stem-and-leaf plot of these page sizes in KB: 41, 47, 52, 38, 55, 49, 63, 44,
  58, 36, 51, 47. Describe the shape in one sentence.
  \item Using the binned table for the 20 latencies (Table~\ref{tab:data-latency-bins}), what
  fraction of requests took under 30 ms? Under 40 ms?
  \item A density histogram of $n=400$ values has a bar over $[2,5)$ of height $0.15$. How many
  values lie in $[2,5)$?
  \item Bins $[0,10)$, $[10,20)$, $[20,50)$ hold 30, 50, 60 values out of $n = 140$. Compute the
  three densities. Which bin has the tallest bar on a density histogram? On a count histogram?
  \item Predict the shape (symmetric, skewed right, skewed left, uniform, bimodal) of: (a) file
  sizes on a laptop; (b) the second-hand reading of a wall clock at random moments; (c) heights
  of a group containing both 6-year-olds and their parents; (d) age at retirement.
  \item Explain why a histogram of 30 values with 30 bins and a histogram with 2 bins would both
  be poor summaries.
  \item In the cache example (Figure~\ref{fig:data-bimodal}), why is ``average latency 37 ms''
  a misleading headline? What two numbers would you report instead?
\end{enumerate}

\subsection*{Answers}

\begin{enumerate}
  \item (a) continuous; (b) discrete; (c) nominal; (d) ordinal; (e) nominal.
  \item $110/250 = 0.44$, $60/250 = 0.24$, $45/250 = 0.18$, $35/250 = 0.14$ (sum 1). A bar chart
  with gaps, bars sorted by height, since language is nominal.
  \item Stems 3 to 6: \texttt{3 | 6 8}, \texttt{4 | 1 4 7 7 9}, \texttt{5 | 1 2 5 8},
  \texttt{6 | 3}. Roughly symmetric, centered in the high 40s, no outliers.
  \item Under 30: $(8+7)/20 = 0.75$. Under 40: $(8+7+2)/20 = 0.85$.
  \item Area $= 0.15 \times 3 = 0.45$, so $0.45 \times 400 = 180$ values.
  \item $30/(140\cdot 10) \approx 0.0214$, $50/(140\cdot 10) \approx 0.0357$,
  $60/(140\cdot 30) \approx 0.0143$. Tallest density bar: $[10,20)$. Tallest count bar:
  $[20,50)$, which misleads because that bin is three times as wide.
  \item (a) skewed right (many small files, a few huge ones); (b) uniform; (c) bimodal;
  (d) roughly symmetric around a typical retirement age, often with a left tail from early
  retirees.
  \item With 30 bins most bins hold 0 or 1 value, so the picture is noise. With 2 bins every
  feature (peaks, skew, outliers) is averaged away.
  \item No request takes anywhere near 37 ms; the data are two groups. Report each group
  separately, e.g.\ typical cache-hit latency (about 20 ms) and typical cache-miss latency
  (about 78 ms), along with the hit rate (105 of 150, 70\%).
\end{enumerate}
